\scanpage{15}

On en conclut : \uline{les structures de similitude} sur le Module (loc. libre de rang 2) $E$ correspondent can. aux sous-Modules loc. monogènes du \inserted{Module (loc. libre de rang 3)} $\underline{\mathrm{End}}(E)/\mathcal O$, i.e. aux sous-Modules $\mathcal A$ de $\underline{\mathrm{End}}(E)$ tels que

\[
\begin{aligned}
1^\circ)&\quad\mathcal A\supset\underline{\mathcal O}\cdot1_E\\
2^\circ)&\quad\mathcal A/\underline{\mathcal O}/1_E \text{ soit loc. monogène.}
\end{aligned}
\]

Les structures de similitudes strictement fidèles correspondent aux sous-fibrés vectoriels de rang 2 (du \inserted{fibré vect. de rang 4} $\underline{\mathrm{End}}(E)$), qui contiennent le sous-fibré $\underline{\mathcal O}\cdot1_E$.

\uline{Lemme 1} Soit $\mathcal A$ un sous-Module de $\underline{\mathrm{End}}(E)$ qui contient $\underline{\mathcal O}\cdot1_E$. Conditions équivalentes

a) $\mathcal A/\underline{\mathcal O}1_E$ est monogène

a$'$) $\exists$ loc. une section $u$ de $\mathcal A$ t.q. $\mathcal A=\underline{\mathcal O}\cdot1_E+\underline{\mathcal O}\cdot u$.

b) $\mathcal A$ est engendré loc. par deux sections.

% slanted
\margin{résulte de l'équation de H.C. $u^2-\mathrm{Tr}(u)\,u+N(u)=0$}

Un tel $\mathcal A$ est stable par \deleted{\unsure{la}} la loi d'algèbre de $\underline{\mathrm{End}}(E)$, et est un faisceau d'anneaux comm. Ainsi \inserted{(op. ou plutôt, \unsure{A de plus})} les $A$ en question \deleted{ainsi} font une \unsure{O|0}). \unsure{revient à dire que}

correspondent : des \inserted{structures de} similitudes sur $E$) sont des sous-algèbres \inserted{commutatives} de $\underline{\mathrm{End}}(E)$. Les str. de similitude str. fidèles correspondent loc\textsuperscript{t} aux $\deleted{\mathcal A}\mathcal A_u$ \deleted{loc. définis par $u$} tels que $u$ ne soit
% archive sheet number: 14
